Pigeonhole Principle: if  n  items in  k  boxes with  n > k,
at least one box has  \ge \lceil n/k \rceil  items.

Pigeonhole (2014 \#3): max holes for 200 pigeons
\frac{k(k+1)}{2} \ge 200 \rightarrow k = 20  (since  \frac{19\cdot20}{2}=190 < 200)
Answer:  k = 19

Exponential equation = 1  (2026 \#14): 
(x^2-5x+5)^{x^2-11x+30} = 1
Case 1: exponent = 0 \rightarrow x^2-11x+30=0 \rightarrow x=5,6
Case 2: base = 1 \rightarrow x^2-5x+4=0 \rightarrow x=1,4
Case 3: base = -1  and exponent even:
x^2-5x+6=0 \rightarrow x=2,3
x=2: exponent=12 (even);\; x=3: exponent=6 (even)
Sum of all solutions: 1+2+3+4+5+6=21

Algebraic identity (2026 \#16):  X = kA+wB,\; A\ne0,\; B\ne0
\frac{X}{AB} - \frac{X-A}{B^2} - \frac{X-B}{A^2} = 0
Substitute and equate coefficients:  k=1,\; w=1 \rightarrow k+w=2

Median/adjusted scores (2026 \#19):
Sorted: 62,68,73,77,78,82,86,86,90,92
Median = \frac{78+82}{2} = 80
Adjustment:  +\frac{80}{10} = +8  to each score
Lowest score:  62 \to 70; percentage increase:  \frac{8}{62} \approx 12.903\%
